\subsection{Exact finite-field computations and their scope}

\begin{claim}{AR.4}{Proved by exhaustive exact computation; the genus-two base-field census}
For $q=4,8,16$, the supplied programs enumerate every
$\mathbf F_q$-isomorphism class of smooth proper genus-two curves,
every rational point of each, and every prime-order constant-deck
pointed cyclic cover over $L=\mathbf F_q$, up to automorphism of its
cyclic deck group.  The exact twisted $L$-polynomial and its relative
class number modulo two are recorded for every test.  The totals are
\[
\begin{array}{c|r|r|r|r}
q&\text{curve classes}&\text{pointed pairs}&\text{tests}&
       \text{pairs with a base-field witness}\\\hline
4&146&730&968&613\\
8&1094&9846&14854&8978\\
16&8470&143990&263765&138769
\end{array}
\]
These are all tests over the displayed base field.  For pairs without
a proved least-prime result, the output makes no assertion of a least
prime over arbitrary extensions.  This is the ``list of all tests
tried'' alternative in A4 for these three genus-two censuses.
\end{claim}
\begin{proof}
We give the completeness argument as well as the reproducible
calculation.  In characteristic two every genus-two curve has a
separable hyperelliptic quotient.  A reduced Artin--Schreier equation
has pole orders $(5)$, $(3,1)$, or $(1,1,1)$, by the different formula
$2g-2=-4+\sum(m_i+1)$.  The first two pole sets are rational; the
last set has split, quadratic, or cubic Frobenius type.  We use
respectively
\[
 ax^5+bx^3+cx+d,\quad ax^3+bx+c/x+d,\quad
 ax+b/x+c/(x+1)+d,
\]
and $ax+(bx+c)/(x^2+x+\tau)+d$ or $n(x)/h(x)+d$ in the last
two cases.  Here $\operatorname{Tr}(\tau)=1$, $h$ is a fixed
irreducible monic cubic, $d\in\{0,\tau\}$, and all required
leading residues are nonzero.  The remaining base transformations
are the full affine group for the first type, scalings for the
second, the permutation group of three branch points in the split
case, the order-two stabilizer of the quadratic set, and the
order-three stabilizer of the cubic set.  A generator change
$y\mapsto y+v(x)$ removes even poles; after reduction it leaves
only the constant Artin--Schreier equivalence.  The program takes
exactly these orbits and does not divide by coefficient Frobenius.

This is the classification in \cite[Propositions~8--16,
Theorem~17]{AR_CNP2005}; as an independent check, our type counts are
\[
\begin{array}{c|rrrrr}
q&f_2=0&f_2=1&\text{split}&\text{quadratic}&\text{cubic}\\\hline
4&14&24&14&48&46\\
8&30&112&140&462&350\\
16&70&480&1240&3930&2750
\end{array}
\]
Their sums agree with that theorem's formula
$2q^3+q^2+q-2+8[4\mid\log_2q]$.  The pole calculations also verify
smoothness of each smooth completion and its genus; no test merely
for rational singularities is used.

For an odd-degree model $Y^2+hY=f$ with unique infinity, we enumerate
all reduced Mumford pairs $(U,V)$, with $U$ monic of degree at most
two, $\deg V<\deg U$, and
$U\mid V^2+hV+f$.  For a fixed $U$ the last condition is an affine
linear system over $\mathbf F_2$ in the coefficients of $V$;
\texttt{finite\_jacobian.py} solves it exactly.  Cantor composition
and reduction implement addition.  Each group element has a unique
such representative, so enumeration is exhaustive.

For an ordinary cubic-pole model with a rational point, we move that
point's image to infinity and translate the Artin--Schreier
coordinate.  The resulting model has $h$ monic of degree three,
$\deg f\le5$, and two rational infinities $O_+,O_-$.
The canonical class is $K=O_++O_-$.  Each nonzero Jacobian class
has a unique representative $E-K$ with $E$ effective of degree two;
the zero class is represented by $K$.  We enumerate
$(U,V,a,b)$ with $\deg U+a+b=2$, allowing infinity multiplicities.
For addition the usual finite-divisor composition first cancels
hyperelliptic pairs.  If both remaining infinity multiplicities
are positive, subtract $K$.  Otherwise the reduction is a
Riemann--Roch calculation in
$\langle1,x,x^2,x^3,Y\rangle=L(3K)$: the code solves the vanishing
conditions on the conjugate finite divisor and the required pole
bounds at $O_+,O_-$.  Dividing the norm of the resulting function
by $U$ gives the complementary divisor.  Its Laurent orders give
the two infinity multiplicities.  This derives the implemented
group law directly from principal divisors.  The script
\texttt{check\_even\_jacobian.py} verifies closure, inverses,
commutativity, and sampled associativity on four independent models.
Pointless models have no pointed tests and are explicitly retained.

For every odd prime $\ell\mid\#J_C(\mathbf F_q)$, the program
forms $J/\ell J$, determines a basis, and enumerates every nonzero
linear functional modulo scalar.  These are exactly all the
surjective prime-order characters, with duplicate generators
removed.  If $v$ is its coordinate vector, we evaluate it on
every $[R-P]$.  AC.4 gives the exact polynomial
\[
 1+\Bigl(\sum_R\zeta_\ell^{\chi(R-P)}\Bigr)T
       +q\zeta_\ell^{\chi(K-2P)}T^2.
\]
Its relative class number is odd precisely when
$1+\sum_R X^{\chi(R-P)}$ is coprime modulo two to
$1+X+\cdots+X^{\ell-1}$, by AC.5.  The program computes that gcd
by binary-polynomial Euclidean division.  Thus every recorded
mod-two value is exact, including primes where two is not a
primitive root.  The JSON records include the quotient basis;
the CSV records give the character coordinates and cyclotomic
polynomial coefficients.  All arithmetic, not only point counting,
is reproduced by the supplied source files and outputs.
\end{proof}

\begin{claim}{AR.5}{Proved by exact computation; degree seven for classes 203 and 302}
Let $\alpha^2+\alpha+1=0$, and let $C_{203}/\mathbf F_4$ have equation
\[
 Y^2+(z^2+z+\alpha)Y
  =(z^2+z+\alpha)(z^3+z^2+z+\alpha^2).
\]
Let $C_{302}$ be its coefficientwise square.  Every original
$\mathbf F_4$-rational point on either curve has an L1 witness of
degree seven over $\mathbf F_4$ or $\mathbf F_{64}$.
Together with the no-degree-three and no-degree-five proof in AC.17,
the least odd prime for each of these fourteen pointed pairs is seven.
The witnesses are constant-deck Lang covers with the indicated
base-point twist, not merely geometric covers.
\end{claim}
\begin{proof}
Both curves have seven rational points and Jacobian order $28$ over
$\mathbf F_4$.  The tests are reproduced by
\texttt{A4\_203\_extensions.py}; all character bases and exact
polynomials are retained in its six JSON outputs.  For $C_{203}$,
the seven points are infinity and the two points above each of
$z=1,\alpha,\alpha^2$.  At infinity the unique prime-seven quotient
has relative class number $9409=97^2$.  At the two points above
$\alpha$ the value is $4243$.  Both values are odd.  At the four
points above $1$ and $\alpha^2$ it is $25768$, so that particular
base-field quotient fails.  Over $\mathbf F_{16}$ the corresponding
tests still fail for these four points.

Over $\mathbf F_{64}$ the Jacobian order is
$6076=4\cdot7^2\cdot31$.  In the deterministically computed basis
of $J/7J$, each of the four remaining points has a successful
character $(1,3)$; the exact norm of its displayed twisted
polynomial is $120974743723$, an odd integer.  The program also
lists the other successful directions and every unsuccessful one.
The infinity direction $(1,0)$ has norm $129436611529$.
Coefficientwise squaring identifies the entire calculation for
$C_{302}$, and the script independently repeats it.

The finite extension comparison here is explicit: each quotient is
of $J_C(\mathbf F_{64})$, and its Lang pullback is based at the
image of the \emph{original} $P$.  Thus it is a geometrically
connected constant cyclic cover in which $P$ splits, and AC.2
applies over that exact field.  We never infer this rational
condition from its holding over a different field.  Finally AC.17
excludes the only smaller odd primes, three and five, over all
finite extensions.  This establishes the asserted minimum, not
only a bound obtained from the search.
\end{proof}

\begin{claim}{AR.6}{Proved by exact computation and Cartier certificates; all F4 pointed pairs}
Every one of the $730$ pointed pairs in the complete $\mathbf F_4$
genus-two census has a rational L1 witness over a specified finite
extension. Starting with the $117$ pairs without a base-field witness,
the successive tests over $\mathbf F_{16},\mathbf F_{64},\mathbf F_{256}$
leave $52,26,24$ pairs. For each of these last $24$, the least prime is
$11$, with a witness over $\mathbf F_{4^{60}}$. An independent
prime-$113$ witness exists over $\mathbf F_{4^{16}}$.
The full characteristic-zero twisted $L$-polynomials are supplied for
the first $706$ pairs. For the last $24$, we give exact mod-two
Cartier certificates and the relative zeta numerator modulo two;
the full characteristic-zero polynomials have not been computed.
This remaining polynomial requirement of A4 is therefore explicit.
\end{claim}
\begin{proof}
The extension scripts use the same exhaustive Jacobian and character
algorithms as AR.4. They embed $\mathbf F_4$ by the first root of
$X^2+X+1$ in the specified larger binary field, map every original
point by that embedding, and base every character at that point.
The irreducible binary moduli are $19,67,283$ for sizes $16,64,256$.
Only pairs already supplied with a witness are omitted at the next
stage. The two final supersingular successes are independently
certified in AC.21 by a degree-five Fermat cover.
The remaining model identifiers are
$42,43,46,47,54,55,56,57$ in the complete F4 JSON census.
In Artin--Schreier form the first four are
\[
 y^2+y=x+c/[x(x+1)]+d,
 \qquad c\in\{\alpha,\alpha^2\},\quad d\in\{0,\alpha\};
\]
the last four are the quadratic-pole forms
$y^2+y=x+c/(x^2+x+\alpha)+d$ with the same parameters.
The first four have respectively $7,3,7,3$ remaining points;
only infinity remains on each of the last four.

All eight forms become one of the two coefficient-conjugate split
forms over $\mathbf F_{4^4}$. Indeed choose $s^2+s=\alpha$, replace
$x$ by $z+s$ in the quadratic-pole forms, and remove the resulting
constant by $y\mapsto y+e$ with $e^2+e=s+d$. These choices exist
in $\mathbf F_{4^4}$. The split twists need only the latter
Artin--Schreier constant change. Both witness fields displayed in
the claim contain this field. The original Jacobian orders
$32,12,48,8$ are prime to $11$ and $113$. Consequently the difference
between any two original rational points is divisible by those
primes already in the original Jacobian, so each corresponding
Lang character has trivial base-point twist on those points.
The coefficient-conjugate form is handled by squaring coefficients.

The exact Cartier calculation, its normalization at infinity,
and the two independent certificates are proved in AC.27--AC.30.
For degree $11$ the relative zeta numerator reduces to
$1+T^4+T^{10}$, whose value at $1$ is one in $\mathbf F_2$.
The degree-$113$ relative numerator is likewise explicitly retained
and has value one. These are proofs of the required rational
relative-class-number parity, rather than point-count estimates.
The no-$3$, no-$5$, and no-$7$ arguments in those statements prove
the asserted minimum $11$ over all finite extensions.

All original finite-field tests, including failures, are retained
in \texttt{A4\_F4\_extensions.csv} and
\texttt{A4\_F4\_F256.csv}. The two large-field certificates retain
an irreducible field modulus, the torsion Mumford divisor, every
nonzero multiple, the tangent coefficients and the Frobenius test.
The independent pure-Python verifier checks these without the
optional compiled arithmetic helper. The combined witness table
identifies these certificates separately from full $L$-polynomial
records; it does not invent the missing characteristic-zero data.
\end{proof}

\begin{claim}{AR.15}{Proved; complete comparison with the geometric R6 classes}
The $146$ arithmetic forms in the $\mathbf F_4$ census map to
exactly $64$ geometric classes: $4$ of two-rank zero, $12$ of
two-rank one, and $48$ ordinary classes. The file
\texttt{A4\_R6\_class\_comparison.csv} gives every geometric class,
every arithmetic model above it, every original rational point's
witness coverage, and the primitive annihilator reported in R6.
The distinguished ordinary classes have the following arithmetic
identifiers:
\[
\begin{array}{c|l}
203&66,67\\
302&64,65\\
122&46,47,56,57\\
133&42,43,54,55.
\end{array}
\]
This verifies the comparison and all original-point rational
certificates. The earlier primitive certificates themselves are
not asserted to have been rerun: their auxiliary coefficient files
were not among the supplied background files.
\end{claim}
\begin{proof}
We compute the geometric invariants, rather than identify arithmetic
models by a suggestive label. For the ordinary normalized equation
$y^2+y=ax+b/x+c/(x+1)$, the unordered triple $\{a,b,c\}$ is
the complete geometric datum. The key $s_1s_2s_3$ records its
elementary symmetric functions. For a quadratic-pole model
$ax+(bx+c)/(x^2+x+\alpha)+d$, choose the two roots $r,r+1$.
Translation gives the triple $(a,br+c,b(r+1)+c)$.
For a cubic-pole model $n(x)/h(x)$, with distinct roots
$r_0,r_1,r_\infty$, put
\[
 K=\frac{r_1+r_\infty}{r_1+r_0},\qquad
 \rho_i=\frac{n(r_i)}{h'(r_i)}.
\]
The substitution $x=(r_\infty z+Kr_0)/(z+K)$ gives
\[
 (a,b,c)=\left(
 \frac{\rho_\infty}{K(r_0+r_\infty)},\quad
 \frac{K\rho_0}{r_\infty+r_0},\quad
 \frac{(1+K)\rho_1}{r_\infty+r_1}\right).
\]
These expressions follow by comparing the simple pole residues at
$z=\infty,0,1$; the remaining term is constant and is removed
geometrically by an Artin--Schreier change.

For two-rank one, $ax^3+bx+c/x$ has normalized pair
$(ac^3,bc)$. For two-rank zero, $ax^5+bx^3+cx$ has invariant
$b^5/a^3$. The completeness of these invariants, including
geometric rather than arithmetic isomorphism, is precisely
\cite[Definition~19, Theorems~20 and~22, Corollary~23]{AR_CNP2005}.
The first two normalizations also follow directly by scaling and
permuting the branch coordinates as above.

The comparison program evaluates all formulas exactly in a common
field $\mathbf F_{4096}$, checks that the symmetric invariants lie
in its specified $\mathbf F_4$ subfield, and obtains the displayed
partition with $64$ keys. It then joins the arithmetic model
identifiers to the independently computed witness table; every
one of the $730$ original points has a successful record.
The R6 column is a transcription of AP-A.C2.22, AP-A.C2.25 and
AP-A.C2.26, labelled as such. In particular its annihilator $15$
for classes $122,133$ is a primitive-factor statement, whereas
the degree-$11$ certificates here are whole-relative rational
statements. Neither column is silently substituted for the other.
\end{proof}

\begin{claim}{AR.16}{Proved; exact cubic classification in all three censuses}
In characteristic two, the locally exact differential theta divisor
of a genus-two curve contains a nonzero point of order three
if and only if the curve is ordinary and its normalized geometric
coefficients satisfy
\[
                         (a+b+c)^3=abc.
\]
Consequently the all-extension cubic-existence criterion AC.22
is decidable from this invariant and the base-field cubic tests.
In the complete pointed censuses, the numbers are
\[
\begin{array}{c|r|r|r}
q&\text{original pairs}&\text{least prime is }3&
   \text{no cubic witness over any extension}\\\hline
4&730&264&466\\
8&9846&3240&6606\\
16&143990&40284&103706.
\end{array}
\]
For $q=4$, explicit characters and full twisted polynomials are
retained for all $264$ positive pairs, over $\mathbf F_4$,
$\mathbf F_{16}$ or $\mathbf F_{256}$. For the other fields the
table is an exact existence and minimum classification; it does
not pretend that every additional character polynomial has been
computed. Whenever a retained degree-five witness has a negative
cubic classification, five is its proved all-extension least prime.
\end{claim}
\begin{proof}
In the ordinary normalized model let $W_0,W_1,W_\infty$ be
the branch points and put $B_\Theta=W_0+W_1-W_\infty$.
The theta divisor consists of $[R-B_\Theta]$. If such a nonzero
class has order three, there is a function with divisor
$3R-3B_\Theta$. Multiplying it by $x(x+1)$ gives a function
with divisor $3R-W_0-W_1-W_\infty$. Riemann--Roch gives
\[
 L(W_0+W_1+W_\infty)=\langle1,y\rangle,
\]
since this divisor has degree three and $g=2$.
Write the function, after scaling, as $y+\beta$, and write
$t=x(R)$ and $\delta=\beta^2+\beta$. The fibre at $R$ is
unramified: a branch point would give a two-primary theta
class, which cannot have order three. Its norm identity is
\[
 b+(b+c)x+ax^2+ax^3+\delta(x^2+x)=a(x+t)^3.
\]
Comparison of coefficients gives
\[
 b=at^3,\qquad c=a(t+1)^3,\qquad \delta=a(t+1).
\]
It follows that $(a+b+c)^3=abc$.

Conversely suppose that equality holds. If $b+c\ne0$, put
$t=(a+c)/(b+c)$; substitution, using the assumed equality,
gives $b=at^3$ and $c=a(t+1)^3$. If $b=c$, the equality forces
$a=b=c$, and either root of $t^2+t+1$ has the same property.
In both cases $t\ne0,1$. Choose $\beta$ with
$\beta^2+\beta=a(t+1)$ and let $R=(t,\beta)$.
The displayed norm identity proves
\[
 \operatorname{div}\left(\frac{y+\beta}{x(x+1)}\right)
                      =3(R-B_\Theta).
\]
The class is nonzero, since $2B_\Theta$ is canonical, whereas
$R$ is not a branch point of the hyperelliptic map. It therefore
has exact order three.

In two-rank zero or one, the locally exact differential line is
$\mathcal O(W)$ for a branch point $W$: in the single-pole model
use $\operatorname{div}(dx)=2W_\infty$, and in the $(3,1)$
model use $\operatorname{div}(dx)=2W_0$.
Thus theta is $C-W$. A nonzero point of order three on theta
would give a function with sole pole of exact order three at
the Weierstrass point $W$. The gap sequence $1,3$, or directly
the basis $L(3W)=\langle1,x\rangle$ after making $W$ infinity,
excludes this. This proves the geometric criterion in every stratum.

AC.22 now says that a cubic witness exists precisely when this
geometric condition holds or a base-field cubic character succeeds.
The latter are all enumerated in AR.4. The program
\texttt{A4\_cubic\_classification.py} computes the invariant
for every arithmetic form and joins that decision to every
original point. Its exact field operations and its full table
give the three counts in the claim. For the ordinary quadratic
model the test is
\[
 (a+b)^3=a(\tau b^2+bc+c^2),\qquad h=x^2+x+\tau;
\]
for cubic-pole models it uses the residue normalization of AR.15.
The required splitting fields have sizes at most $q^3$.

For the $59$ positive $\mathbf F_4$ pairs lacking a base-field
cubic witness, the extra exact character tests supply $49$
over $\mathbf F_{16}$ and the remaining $10$ over
$\mathbf F_{256}$. Both successful and unsuccessful tests are
retained. In particular the latter ten are not assigned the
incorrect smaller field merely from the geometric theta condition.
Finally there is no smaller admissible prime than three; after
three is excluded, a successful degree-five record proves its
minimum as stated.
\end{proof}

\begin{claim}{AR.7}{Open; the remaining computational scope of A4}
The base-field pointed odd-prime tests for smooth plane quartics over
$\mathbf F_3$ and $\mathbf F_5$ are complete in AG.24 and AG.26,
including independent exact integer norm audits. The complete
consolidated table is \path{code/A4_quartic_witnesses.csv}.
However, least primes over arbitrary finite extensions have not been
determined for every genus-two or quartic pointed pair outside the
explicitly proved cases. The full characteristic-zero character
polynomials for the large-field degree-$11$ and degree-$113$
genus-two certificates have not been computed; their full mod-two
certificates and the precise fields are retained in AC.28.
The exhaustive tables and successful finite tests do not replace
either missing minimum or missing characteristic-zero polynomial.
\end{claim}
